% !TeX program = xelatex
\documentclass[12pt,a4paper]{article}

\usepackage{fontspec}
\usepackage[russian]{babel}
\setmainfont{DejaVu Serif}
\setsansfont{DejaVu Sans}
\setmonofont{DejaVu Sans Mono}
\usepackage{geometry}
\usepackage{indentfirst}
\usepackage{setspace}
\usepackage{enumitem}
\usepackage{array}
\usepackage{longtable}
\usepackage{tabularx}
\usepackage{booktabs}
\usepackage{graphicx}
\usepackage{float}
\usepackage{xcolor}
\usepackage{tikz}
\usepackage[hidelinks]{hyperref}

\usetikzlibrary{arrows.meta,positioning,shapes.geometric,fit,calc}
\geometry{left=2.5cm,right=1.5cm,top=2cm,bottom=2cm}
\onehalfspacing
\setlength{\parindent}{1.25cm}
\setlength{\parskip}{0pt}
\setlist{nosep,leftmargin=1.25cm}
\renewcommand{\arraystretch}{1.25}
\sloppy

\definecolor{itmoblue}{RGB}{27,74,132}
\definecolor{lightblue}{RGB}{232,240,249}
\definecolor{lightgray}{RGB}{245,245,245}

\newcommand{\systemname}{«ЛабРезерв»}
\newcommand{\req}[1]{\textbf{#1}}
\newcommand{\uc}[1]{\textbf{#1}}

% Сценарии прецедентов оформляются построчно: каждый шаг визуально отделён.
\newlist{ucsteps}{enumerate}{1}
\setlist[ucsteps]{label=\arabic*),leftmargin=*,itemsep=5pt,topsep=4pt,parsep=0pt}
\newlist{ucalts}{enumerate}{1}
\setlist[ucalts]{label=A\arabic*:,leftmargin=*,itemsep=5pt,topsep=4pt,parsep=0pt}


\usepackage{pdfpages}
\usepackage{fvextra}
\usepackage{placeins}
\usepackage{needspace}
\setlength{\emergencystretch}{2em}
\newcommand{\code}[1]{\texttt{\detokenize{#1}}}
\makeatletter
\renewcommand*\l@section[2]{\@dottedtocline{1}{0em}{2.5em}{\bfseries #1}{\bfseries #2}}
\renewcommand*\l@subsection{\@dottedtocline{2}{1.5em}{3.5em}}
\makeatother
\begin{document}
\hypersetup{pageanchor=false}
\pagenumbering{gobble}

\begin{titlepage}
\begin{center}
    \large
    Национальный исследовательский университет ИТМО\\
    Факультет программной инженерии и компьютерной техники\\
    \vspace{2.5cm}
    \textbf{\Large Информационные системы}\\
    \vspace{0.5cm}
    \textbf{\Large Курсовая работа}\\
    \vspace{0.35cm}
    \textbf{\large Этап 2. Проектирование и реализация\\базы данных}\\
    \vspace{0.8cm}
    \textbf{\Large Система управления бронированием\\
    и загрузкой оборудования в лаборатории/цеху}\\
    \vspace{2.4cm}
    \begin{flushright}
        Выполнили студенты 3 курса:\\
        Калинин Дмитрий Викторович\\
        Рзаев Ахмед Меджидович\\
        \vspace{0.6cm}
        Преподаватель:\\
        Николаев Владимир Вячеславович
    \end{flushright}
    \vfill
    Санкт-Петербург, 2026
\end{center}
\end{titlepage}

\clearpage
\hypersetup{pageanchor=true}
\pagenumbering{roman}
\begin{singlespace}
\tableofcontents
\end{singlespace}
\newpage
\pagenumbering{arabic}

\section{Введение и статус работы}

Цель второго этапа --- преобразовать утверждённые предметную область, требования
и прецеденты системы \systemname{} в воспроизводимую базу PostgreSQL.
Результат включает концептуальную ER-модель, даталогическую схему, SQL-скрипты,
ограничения, триггеры, критические функции, индексы и проверочные сценарии.
Первый этап используется как неизменяемая основа: \code{stage-01/report_approved.pdf}
и соответствующий \code{report.tex}.

\textbf{Статус ER-модели: «Ожидает согласования».} Преподаватель ещё не подтвердил
эту модель. Согласно официальному заданию согласование предшествует
даталогической модели; по поручению авторов остальные материалы реализованы
на рабочей версии, чтобы предъявить законченный черновик. После согласования
правки вносятся в исходники и SQL, затем проверки выполняются заново.

Java-приложение пока не реализовано. Будущий модульный монолит
\code{labreserve.war} развёртывается во внешнем WildFly на Helios;
прикладной доступ выполняется через Spring Data JPA/Hibernate.
Критические функции PostgreSQL второго этапа должны использоваться на третьем
этапе через нативные JPA-запросы. Spring JDBC не используется.

\section{Анализ требований к данным}

Модель сохраняет три сквозных процесса первого этапа, не подменяя их набором CRUD.
Мелкие действия остаются декомпозицией UC-01--UC-03, представленной через
\code{include}/\code{extend} в первом отчёте.

\begin{enumerate}
\item UC-01: каталог и допуски, календарь, атомарное подтверждение, конфликт,
альтернативы, выбор, перенос, отмена, фактическая работа.
\item UC-02: сообщение о неисправности или ТО, проверка Техником либо
Администратором, недоступность, затронутые брони, варианты и уведомления,
ремонт и восстановление доступности.
\item UC-03: типы, функции, единицы оборудования, роли, допуски, история,
плановая и фактическая загрузка, простои, отказы и выгрузка аналитики.
\end{enumerate}

Основные инварианты утверждены как BR-1--BR-10. Допуск обязателен,
подтверждённые интервалы не пересекаются, недоступность блокирует новые брони,
окончание строго позже начала. Альтернативы учитывают функцию и приоритеты.
Будущие затронутые брони требуют переноса, но выбор не делается автоматически.
Терминальные брони и история сохраняются; один фактический сеанс закрывается
единожды. Временные данные учитывают часовые пояса.

FR-01 и NFR безопасности имеют границу ответственности: БД хранит роли и
проверяет переданную идентичность, а аутентификацию, HTTPS и безопасное
хеширование паролей реализует будущий сервер. Матрица ниже указывает эту
границу, а не объявляет весь будущий продукт реализованным.

\subsection{Матрица трассируемости}
В каждой строке указаны объекты схемы, операция и проверка. T01--T60 ---
утверждающие SQL-тесты; C01--C05 --- конкурентные сценарии; Q1--Q5 --- планы
запросов. Полная машинно читаемая и Markdown-версия сохранена в
\code{docs/TRACEABILITY.md}.
\begin{singlespace}\small
\begin{longtable}{>{\raggedright\arraybackslash}p{0.1\textwidth}>{\raggedright\arraybackslash}p{0.24\textwidth}>{\raggedright\arraybackslash}p{0.26\textwidth}>{\raggedright\arraybackslash}p{0.3\textwidth}}
\toprule Код & Объекты & Операция & Проверка\\\midrule\endfirsthead
\toprule Код & Объекты & Операция & Проверка\\\midrule\endhead
UC-01 & booking, admission, alternative, use\_\allowbreak{}session & create/\allowbreak{}select/\allowbreak{}open/\allowbreak{}close & T01–11,16–19,30–34,44,52–53; C01,C04 \\
UC-02 & unavailability, fault\_\allowbreak{}report, repair, notification & register/\allowbreak{}edit/\allowbreak{}resolve; impact & T12–29; C02,C03,C05 \\
UC-03 & catalogue, admission, audit, conflict & grant/\allowbreak{}revoke; metrics & T35–40,45–50,56; Q3–Q5 \\
FR-01 & app\_\allowbreak{}user, user\_\allowbreak{}role & assert\_\allowbreak{}owner, require\_\allowbreak{}staff & T11,25,45,58–60; аутентификация — этап 3 \\
FR-02 & equipment\_\allowbreak{}type, equipment, type\_\allowbreak{}capability & DDL, audit\_\allowbreak{}change & seed; интерфейс администратора — этап 3 \\
FR-03 & equipment, unavailability & equipment\_\allowbreak{}state & T22,27,29 \\
FR-04 & admission & grant/\allowbreak{}revoke\_\allowbreak{}admission & T45–50 \\
FR-05 & equipment, booking, unavailability & free\_\allowbreak{}intervals & T42; Q1–Q3 \\
FR-06 & booking, capability & create\_\allowbreak{}booking & T01,09 \\
FR-07 & admission, booking, unavailability & check\_\allowbreak{}booking & T02–04,15,24,47–49 \\
FR-08 & booking, equipment & create\_\allowbreak{}booking; EXCLUDE & C01,C04,C05 \\
FR-09 & booking, notification & cancel\_\allowbreak{}booking & T06,07,11 \\
FR-10 & booking, unavailability & free\_\allowbreak{}intervals, find\_\allowbreak{}alternatives & T42,46 \\
FR-11 & type\_\allowbreak{}capability, admission & equivalent\_\allowbreak{}equipment & T40 \\
FR-12 & alternative & find/\allowbreak{}offer\_\allowbreak{}alternatives & T13,19,26,28 \\
FR-13 & booking, alternative & select\_\allowbreak{}alternative & T16–18,44 \\
FR-14 & unavailability & register/\allowbreak{}edit\_\allowbreak{}maintenance & T12,20 \\
FR-15 & booking, alternative & impact\_\allowbreak{}bookings & T12–15 \\
FR-16 & fault\_\allowbreak{}report, unavailability & report\_\allowbreak{}fault, register\_\allowbreak{}unavailability & T21–25 \\
FR-17 & alternative, notification & impact\_\allowbreak{}bookings & T23,26,28 \\
FR-18 & repair, unavailability & resolve\_\allowbreak{}unavailability & T27,29 \\
FR-19 & use\_\allowbreak{}session, booking & open/\allowbreak{}close\_\allowbreak{}session & T30–34,52,53 \\
FR-20 & booking, use\_\allowbreak{}session & equipment\_\allowbreak{}metrics & T57 \\
FR-21 & notification & cancel/\allowbreak{}select/\allowbreak{}impact/\allowbreak{}reminders & T14,43 \\
FR-22 & booking, use\_\allowbreak{}session, conflict\_\allowbreak{}event & equipment\_\allowbreak{}metrics & T38,39,56 \\
FR-23 & equipment\_\allowbreak{}metrics result & export\_\allowbreak{}report.sql & CSV выполнен; UI-фильтры — этап 3/\allowbreak{}4 \\
FR-24 & audit\_\allowbreak{}event & audit\_\allowbreak{}change; audit\_\allowbreak{}immutable & T35,36 \\
BR-1 & admission, app\_\allowbreak{}user & check\_\allowbreak{}booking; grant/\allowbreak{}revoke & T02,03,47–49 \\
BR-2 & ex\_\allowbreak{}booking\_\allowbreak{}overlap & EXCLUDE + lock\_\allowbreak{}equipment & T04,05; C01,C04 \\
BR-3 & unavailability & booking\_\allowbreak{}guard; impact & T15,24,27; C02,C03 \\
BR-4 & time\_\allowbreak{}window & CHECK + диапазон & T08,54 \\
BR-5 & type\_\allowbreak{}capability & equivalent/\allowbreak{}find; check\_\allowbreak{}booking & T09,40 \\
BR-6 & alternative.priority & find\_\allowbreak{}alternatives & T19 \\
BR-7 & booking.status & impact; select\_\allowbreak{}alternative & T12,20,23 \\
BR-8 & audit, booking, admission & history guards & T07,34–36,41,50,51 \\
BR-9 & use\_\allowbreak{}session & session guards; open/\allowbreak{}close & T10,30–33,52,53 \\
BR-10 & timestamptz, tstzrange & time\_\allowbreak{}window; TZ отображения & T05,37 \\
\bottomrule\end{longtable}\end{singlespace}

\section{Концептуальная ER-модель}

Выделены \textbf{15 содержательных концептуальных сущностей}.
При преобразовании M:N вводятся ещё две связующие таблицы, поэтому
даталогическая модель содержит 17 таблиц. Допуск --- отдельная ассоциативная
сущность с периодом, основанием, автором и отзывом, а не простой список пар.

Две явно моделируемые M:N:
\begin{enumerate}
\item Пользователь --- Роль. Человек может совмещать роли; одна роль назначается
многим людям. Реляционная таблица \code{user_role} имеет составной PK.
\item Тип оборудования --- Функциональная возможность. Тип поддерживает
несколько функций; одну функцию поддерживают разные типы. Например, раскрой
листа возможен лазером и подходящим ЧПУ. Таблица \code{type_capability} устраняет
дублирование списка функций у каждой единицы оборудования.
\end{enumerate}

Модель представлена тремя тематическими листами с общей идентичностью повторённых
сущностей. Исходники \code{diagrams/conceptual.tex} (TikZ) и
\code{diagrams/conceptual.mmd} (Mermaid) редактируемы.
Графические листы показывают основные связи; дополнительные ссылки на авторов
указаны в подписях и полностью перечислены в даталогической схеме.

\subsection{Кратности и обязательность}
\begin{itemize}
\item Каждый обязательный FK дочерней записи соответствует ровно одному
родителю; родитель может иметь 0..N детей.
\item Бронь имеет 0..1 фактический сеанс; UNIQUE по booking\_id запрещает второй.
\item Сообщение имеет 0..1 подтверждённый период неисправности. Для поломки
ссылка на сообщение обязательна, для ТО её не должно быть.
\item Получатель уведомления обязателен; связь уведомления с бронью необязательна.
\item Автор аудита необязателен для миграционных и системных действий.
Объект аудита связан логически через таблицу и ключ, без полиморфного FK.
\item Автор регистрации недоступности обязателен; автор завершения появляется
только при закрытии. Исполнитель технической работы обязателен.
\end{itemize}

\clearpage
\includepdf[pages=-,fitpaper=true,pagecommand={\thispagestyle{empty}\begin{tikzpicture}[remember picture,overlay]\node[anchor=south] at ([yshift=5mm]current page.south){\thepage};\end{tikzpicture}}]{diagrams/conceptual.pdf}

\section{Преобразование в даталогическую модель}

Использована схема \code{labreserve}; названия таблиц и столбцов ---
\code{snake_case}, индексы --- \code{ix_...}, уникальные индексы --- \code{ux_...},
существенные ограничения --- \code{ck_...}/\code{ex_...}.
Остальные имена PK/FK/CHECK/UNIQUE назначаются единообразно PostgreSQL.
Все сущности с самостоятельной идентичностью используют
\code{bigint GENERATED ALWAYS AS IDENTITY}. Связующие таблицы имеют составные PK;
роль идентифицируется фиксированным кодом.

\subsection{Таблицы и назначение}
\begin{singlespace}\small
\begin{longtable}{>{\raggedright\arraybackslash}p{0.32\textwidth}>{\raggedright\arraybackslash}p{0.6\textwidth}}
\toprule Таблица & Назначение\\\midrule\endfirsthead
\toprule Таблица & Назначение\\\midrule\endhead
app\_\allowbreak{}user & Учётная запись человека; активность и профиль \\
app\_\allowbreak{}role & Три внешние роли предметной области \\
user\_\allowbreak{}role & Назначение ролей пользователям (M:N) \\
equipment\_\allowbreak{}type & Тип оборудования и архивирование \\
capability & Справочник требуемых функций \\
type\_\allowbreak{}capability & Поддерживаемые типом функции (M:N) \\
equipment & Инвентарная единица, характеристики, архив \\
admission & Допуск к типу; период, основание, автор и отзыв \\
booking & Плановый интервал, владелец, цель и жизненный цикл \\
alternative & Историческое предложение переноса и выбор \\
fault\_\allowbreak{}report & Первичное сообщение и результат проверки \\
unavailability & Период ТО либо подтверждённой неисправности \\
repair & Исполнитель и результат технической работы \\
use\_\allowbreak{}session & Фактические начало и окончание работы \\
notification & Внутрисистемное уведомление и прочтение \\
audit\_\allowbreak{}event & Неизменяемый журнал значимых изменений \\
conflict\_\allowbreak{}event & Отклонённый запрос для статистики конфликтов \\
\bottomrule\end{longtable}\end{singlespace}

\subsection{Нормализация}
Справочные названия не копируются в брони; роль и функции вынесены в M:N,
сеанс отделён от планового интервала. Реляционные атрибуты зависят от ключа;
схема ориентирована на 3НФ. Обоснованные специальные случаи:
\begin{enumerate}
\item JSONB характеристик позволяет хранить разные наборы физических параметров;
он не заменяет нормализованный справочник функций и не определяет эквивалентность.
\item JSONB аудита --- исторический снимок до/после изменения, а не основной
источник текущих данных.
\item Интервал и приоритет альтернативы сохраняются как снимок предложения;
перед выбором их актуальность проверяется повторно.
\item schedule\_revision --- технический счётчик сериализации, не характеристика
оборудования. Изменения только счётчика не добавляются в аудит.
\end{enumerate}

Техническое состояние вычисляется функцией \code{equipment_state}:
архивирование выше неисправности, неисправность выше ТО, иначе оборудование
доступно. Оно не хранится вторым расходящимся значением.

\subsection{Типы времени и ключи}
Моменты времени имеют тип \code{timestamptz}, интервалы --- домен
\code{time_window} на основе \code{tstzrange}. Границы полуоткрытые
$[start,end)$; соседние интервалы допустимы. Домен запрещает пустой интервал,
неограниченное начало и включённый конец. Брони, варианты, ТО и технические
работы имеют конечный конец. Для поломки конец может быть неизвестен до ремонта.

\code{timestamptz} сохраняет момент, а не исходное написание часового пояса.
При отображении применяется зона лаборатории \code{Europe/Moscow}; два ввода
\code{10:00+03} и \code{07:00Z} обозначают одинаковое время.
Все 27 FK имеют ON DELETE RESTRICT и ON UPDATE RESTRICT. Идентификаторы
не изменяются; вместо удаления используется архивирование, отзыв, отмена или
закрытие. Полный перечень полей, NULL, default и ограничений дан в приложении А.

\clearpage
\includepdf[pages=-,fitpaper=true,pagecommand={\thispagestyle{empty}\begin{tikzpicture}[remember picture,overlay]\node[anchor=south] at ([yshift=5mm]current page.south){\thepage};\end{tikzpicture}}]{diagrams/logical.pdf}

\section{Обеспечение целостности и история}

\subsection{Декларативные средства}
PK и UNIQUE обеспечивают идентичность, уникальность логинов и инвентарных номеров,
связующих пар, сеанса брони и подтверждения сообщения. FK запрещают сирот.
CHECK контролируют перечисления статусов, метаданные отмены и завершения,
JSONB-объект характеристик и корректные интервалы.

Частичное ограничение \code{ex_booking_overlap} проверяет одновременно
равенство оборудования и пересечение диапазонов. GiST-класс для bigint
предоставляет расширение \code{btree_gist}. В предикат включён только
\code{status='confirmed'}: отменённые, завершённые и требующие переноса брони
не занимают плановое расписание.

Ограничение одной таблицы \code{booking} не проверяет отдельную
\code{unavailability}. Межтабличная защита требует единого протокола блокировки
и триггеров; обычный SELECT перед INSERT недостаточен.

\subsection{Протокол конкурентных операций}
\begin{enumerate}
\item Бронь, технический период и сеанс используют одну строку оборудования.
\item \code{lock_equipment} обновляет счётчик строки; блокировка удерживается
до конца внешней транзакции.
\item При READ COMMITTED конкурентная операция ждёт; следующие SQL-команды
VOLATILE-функции читают уже зафиксированные данные.
\item При REPEATABLE READ/Serializable изменение строки после снимка приводит
к 40001. Проверка на устаревшем снимке не продолжается.
\item После регистрации недоступности будущие пересекающиеся подтверждённые
брони получают needs\_transfer. Уведомления и предложения создаются атомарно.
\item При переносе обе единицы блокируются в порядке id; триггер повторяет
проверки допуска, функции, архива и технических периодов, а EXCLUDE проверяет брони.
\item Прямой UPDATE может получить блокировку брони раньше оборудования и
вызвать 40P01 при обратном порядке блокировок. Это безопасный отказ; сервер
должен повторить всю транзакцию при 40001/40P01 ограниченное число раз.
\end{enumerate}

\subsection{Триггеры и аудит}
В схеме \textbf{25 пользовательских триггеров}:
\begin{itemize}
\item booking\_guard: активность, допуск на весь интервал, функция, архив,
технические периоды, неизменность владельца и терминальных броней.
\item unavailability\_guard: Техник или Администратор, соответствующее сообщение,
неизменность завершённой причины.
\item unavailability\_impact: будущие брони, варианты и уведомления.
\item session\_guard и session\_complete: начало по confirmed, актуальный допуск
и состояние, попадание в плановый интервал, одно закрытие, completed.
\item repair\_guard, fault\_guard: исполнители и история технической проверки.
\item audit\_immutable, conflict\_immutable, repair\_immutable, fault\_no\_delete,
admission\_guard: сохранение исторических записей.
\item 13 audit\_*: значимые изменения справочников и операций, старое/новое
значение, автор, время и идентификатор транзакции. Хеши паролей в снимки не попадают.
\end{itemize}

Невозможно произвольно переписать completed/cancelled. Исправление описывается
отдельной записью аудита. Завершённые сеансы и ремонты неизменяемы.
Отзыв допуска сохраняет старую запись; новый срок требует нового допуска.
Закрытие одной причины недоступности не отменяет другие причины.

\subsection{Доверие и права}
Владелец схемы выполняет миграции. PUBLIC не имеет прав на схему, таблицы и функции.
\code{grant_application.sql} выдаёт существующей служебной роли backend SELECT
и EXECUTE только внешнего API. Прямого DML расписания нет.
Изменяющие внешние функции имеют SECURITY DEFINER и фиксированный search\_path;
внутренние функции не выдаются backend. T58--T60 проверяют эти права реально.

Идентичность человека удостоверяет будущий Spring Security. Значение p\_actor
должно приходить из серверного контекста; оно не является доказательством
личности при прямом доступе к БД. Веб-клиенты к PostgreSQL не подключаются.
Владелец/суперпользователь может отключать триггеры; это административная граница,
а не обещание защиты от владельца БД.

\section{SQL-скрипты и демонстрационные данные}

Скрипт \code{install.sql} выполняет четыре части в общей транзакции: таблицы и
домен, триггеры, операции, индексы. При ошибке установка откатывается.
Отдельные скрипты создают и удаляют только \code{labreserve_stage02_test}.
Очистка схемы дополнительно проверяет имя текущей БД; DROP DATABASE не использует
FORCE. Скрипты не удаляют реальные пользовательские базы.

Учебные данные включают:
\begin{enumerate}
\item Пять людей: пользователь, администратор, техник, новичок без допуска,
неактивный пользователь. Администратор совмещает две роли.
\item Четыре типа и четыре функции. Раскрой доступен двум типам, ЧПУ имеет две функции.
\item Пять единиц: два 3D-принтера, лазерный резак, ЧПУ, измерительный прибор.
\item Двенадцать допусков с основаниями и авторами.
\item Требующие переноса, отменённая и завершённая брони; предложения, ТО,
сообщения, активная поломка, завершённый ремонт, фактический сеанс и уведомления.
\end{enumerate}

Времена демонстрации относительны к дню запуска. Хеши паролей не подставлены:
аутентификация демонстрационных людей пока отсутствует. Начальные идентификаторы
предполагают пустую схему; seed намеренно не идемпотентен. Runner пересоздаёт
только выделенную учебную схему.

Полный запуск и сборка описаны в \code{README.md}. Fedora требует пакеты
PostgreSQL server/contrib, Python 3, XeLaTeX, Poppler и шрифты DejaVu.
CSV формируется \code{export_report.sql}; \code{export_model.sql} читает
системные каталоги и сохраняет фактическую модель в JSON.

\section{Критические функции PostgreSQL}

Внешний SQL API включает \textbf{19 функций}; ещё 13 функций обеспечивают
внутренние проверки и триггерные обработчики. Читающие операции используют SQL
либо STABLE PL/pgSQL; изменяющие критические операции --- PL/pgSQL.
Отдельные процедуры не требуются: возврат идентификатора/результата функции удобен
для следующего слоя JPA; COMMIT выполняет внешняя транзакция, не сама функция.

Общие ошибки: P0002 --- объект не найден; 23503 --- неверный FK;
23514 --- бизнес-ограничение; 23P01 --- конфликт брони; 42501 --- нет полномочий.
40001/40P01 требуют повтора всей транзакции. Остальные ошибки не маскируются.
Контракты параметров, результата, ошибок и покрытия перечислены ниже.
\subsection{free\_\allowbreak{}intervals}
Максимальные свободные диапазоны в конечном окне; можно исключить переносимую бронь.
\begin{itemize}
\item Вход: p\_\allowbreak{}equipment bigint,p\_\allowbreak{}window tstzrange,p\_\allowbreak{}ignore bigint DEFAULT NULL.
\item Выход: SETOF tstzrange.
\item Ошибки: 22023: некорректное окно.
\item Гарантии: Предварительный снимок; объединение броней и технических периодов.
\item Трассировка: UC-01; FR-05,10; BR-2,3,10; T42.
\end{itemize}
\subsection{equivalent\_\allowbreak{}equipment}
Свободные единицы с требуемой функцией и допуском на весь интервал.
\begin{itemize}
\item Вход: p\_\allowbreak{}user bigint,p\_\allowbreak{}capability bigint,p\_\allowbreak{}period tstzrange.
\item Выход: SETOF equipment.
\item Ошибки: Пустое множество при отсутствии вариантов.
\item Гарантии: Не резервирует найденные слоты.
\item Трассировка: UC-01; FR-11; BR-1,5; T40.
\end{itemize}
\subsection{find\_\allowbreak{}alternatives}
До 12 вариантов в горизонте; пересекает свободные интервалы со сроками допусков.
\begin{itemize}
\item Вход: p\_\allowbreak{}user bigint,p\_\allowbreak{}equipment bigint,p\_\allowbreak{}capability bigint,  p\_\allowbreak{}period tstzrange,p\_\allowbreak{}horizon interval DEFAULT interval '30 days',p\_\allowbreak{}ignore bigint DEFAULT NULL.
\item Выход: equipment\_\allowbreak{}id, period, priority.
\item Ошибки: 22023: неверный диапазон или горизонт.
\item Гарантии: Сортировка BR-6: 1,2,3; затем время и id.
\item Трассировка: UC-01/\allowbreak{}02; FR-10–12; BR-5,6; T13,19,26,28,46.
\end{itemize}
\subsection{offer\_\allowbreak{}alternatives}
Сохраняет предложения к существующей брони; прежние предложения устаревают.
\begin{itemize}
\item Вход: p\_\allowbreak{}booking bigint,p\_\allowbreak{}horizon interval DEFAULT interval '30 days'.
\item Выход: integer: число предложений.
\item Ошибки: P0002: бронь отсутствует.
\item Гарантии: Выбор пользователя не имитируется; предложения не занимают слот.
\item Трассировка: UC-01/\allowbreak{}02; FR-12,15,17; BR-7; T13,20.
\end{itemize}
\subsection{create\_\allowbreak{}booking}
Подтверждает будущую бронь; p\_\allowbreak{}user — удостоверенный пользователь.
\begin{itemize}
\item Вход: p\_\allowbreak{}user bigint,p\_\allowbreak{}equipment bigint,p\_\allowbreak{}capability bigint,  p\_\allowbreak{}start timestamptz,p\_\allowbreak{}end timestamptz,p\_\allowbreak{}purpose text.
\item Выход: bigint: id брони.
\item Ошибки: 23514: допуск/\allowbreak{}архив/\allowbreak{}недоступность; 23P01: пересечение; 22000: обратный диапазон.
\item Гарантии: Атомарный INSERT, триггер и EXCLUDE.
\item Трассировка: UC-01; FR-06–08; BR-1–4; T01–09; C01,C04.
\end{itemize}
\subsection{request\_\allowbreak{}booking}
Обёртка создания: сохраняет причину отклонения без исключения для ожидаемого конфликта.
\begin{itemize}
\item Вход: p\_\allowbreak{}user bigint,p\_\allowbreak{}equipment bigint,p\_\allowbreak{}capability bigint,  p\_\allowbreak{}start timestamptz,p\_\allowbreak{}end timestamptz,p\_\allowbreak{}purpose text.
\item Выход: booking\_\allowbreak{}id либо error\_\allowbreak{}code.
\item Ошибки: Ожидаемые 23514/\allowbreak{}23P01 возвращаются; прочие исключения передаются выше.
\item Гарантии: Запись отказа сохранится только при COMMIT внешней транзакции.
\item Трассировка: UC-01/\allowbreak{}03; FR-08,22; T38.
\end{itemize}
\subsection{cancel\_\allowbreak{}booking}
Владелец либо администратор отменяет будущую активную бронь с причиной.
\begin{itemize}
\item Вход: p\_\allowbreak{}actor bigint,p\_\allowbreak{}booking bigint,p\_\allowbreak{}reason text.
\item Выход: void.
\item Ошибки: 42501: чужая бронь; 23514: прошлая/\allowbreak{}терминальная бронь или существующий сеанс.
\item Гарантии: Отмена и уведомление атомарны; история сохраняется.
\item Трассировка: UC-01; FR-09,21; BR-8; T06,07,11.
\end{itemize}
\subsection{select\_\allowbreak{}alternative}
Владелец либо администратор подтверждает сохранённое предложение.
\begin{itemize}
\item Вход: p\_\allowbreak{}actor bigint,p\_\allowbreak{}alternative bigint.
\item Выход: void.
\item Ошибки: 23514: устаревший вариант; 23P01: конкурентная бронь; 42501: нет прав.
\item Гарантии: Блокировка обеих единиц; повторная проверка допуска, функции и календаря.
\item Трассировка: UC-01/\allowbreak{}02; FR-13; BR-5–7; T16–18,44.
\end{itemize}
\subsection{report\_\allowbreak{}fault}
Активный человек сообщает о неисправности без автоматического изменения доступности.
\begin{itemize}
\item Вход: p\_\allowbreak{}actor bigint,p\_\allowbreak{}equipment bigint,p\_\allowbreak{}description text.
\item Выход: bigint: id сообщения.
\item Ошибки: 42501: неактивный автор; 23503: нет оборудования.
\item Гарантии: Сохраняется сообщение и аудит.
\item Трассировка: UC-02; FR-16; T21.
\end{itemize}
\subsection{register\_\allowbreak{}unavailability}
Техник либо администратор регистрирует ТО или подтверждает поломку. NULL p\_\allowbreak{}end — неизвестный конец поломки.
\begin{itemize}
\item Вход: p\_\allowbreak{}actor bigint,p\_\allowbreak{}equipment bigint,p\_\allowbreak{}kind text,  p\_\allowbreak{}start timestamptz,p\_\allowbreak{}end timestamptz,p\_\allowbreak{}reason text,p\_\allowbreak{}report bigint DEFAULT NULL.
\item Выход: bigint: id периода.
\item Ошибки: 42501: роль; 23514: несоответствие сообщения/\allowbreak{}периода; 23503: ссылки.
\item Гарантии: Период, маркировка будущих броней, варианты и уведомления в одной транзакции.
\item Трассировка: UC-02; FR-14–17; BR-3,7; T12–15,22–25; C02,C03,C05.
\end{itemize}
\subsection{edit\_\allowbreak{}maintenance}
Изменение периода и причины активного ТО либо его отмена.
\begin{itemize}
\item Вход: p\_\allowbreak{}actor bigint,p\_\allowbreak{}unavailability bigint,p\_\allowbreak{}start timestamptz,  p\_\allowbreak{}end timestamptz,p\_\allowbreak{}reason text,p\_\allowbreak{}cancel boolean DEFAULT false.
\item Выход: void.
\item Ошибки: 23514: неактивное ТО/\allowbreak{}иная причина; 42501: роль.
\item Гарантии: Новые затронутые брони маркируются; отмена не подтверждает их сама.
\item Трассировка: UC-02; FR-14,15; BR-7,8; T20.
\end{itemize}
\subsection{equipment\_\allowbreak{}state}
Состояние на момент p\_\allowbreak{}at: archived, fault, maintenance, available.
\begin{itemize}
\item Вход: p\_\allowbreak{}equipment bigint,p\_\allowbreak{}at timestamptz DEFAULT clock\_\allowbreak{}timestamp().
\item Выход: text.
\item Ошибки: NULL при неизвестной единице.
\item Гарантии: Единый вычисляемый источник состояния по всем причинам.
\item Трассировка: UC-02/\allowbreak{}03; FR-03; T22,27,29.
\end{itemize}
\subsection{resolve\_\allowbreak{}unavailability}
Техник либо администратор записывает результат и закрывает активную причину.
\begin{itemize}
\item Вход: p\_\allowbreak{}actor bigint,p\_\allowbreak{}unavailability bigint,p\_\allowbreak{}result text.
\item Выход: text: итоговое состояние.
\item Ошибки: 23514: причина уже закрыта/\allowbreak{}ещё не началась; 42501: роль.
\item Гарантии: Сохраняется вид работы work\_\allowbreak{}kind (maintenance/\allowbreak{}repair); остальные события продолжают блокировать.
\item Трассировка: UC-02; FR-18; BR-3,8; T27,29.
\end{itemize}
\subsection{open\_\allowbreak{}session}
Владелец либо администратор открывает фактическую работу в плановом интервале.
\begin{itemize}
\item Вход: p\_\allowbreak{}actor bigint,p\_\allowbreak{}booking bigint.
\item Выход: bigint: id сеанса.
\item Ошибки: 23514: бронь/\allowbreak{}допуск/\allowbreak{}время/\allowbreak{}доступность; 23505: сеанс уже есть.
\item Гарантии: Один сеанс; блокировка оборудования и проверка начала.
\item Трассировка: UC-01; FR-19,20; BR-9; T10,30,31.
\end{itemize}
\subsection{close\_\allowbreak{}session}
Закрывает работу. p\_\allowbreak{}comment обязателен при действии администратора за владельца.
\begin{itemize}
\item Вход: p\_\allowbreak{}actor bigint,p\_\allowbreak{}session bigint,p\_\allowbreak{}comment text DEFAULT NULL.
\item Выход: void.
\item Ошибки: 23514: повторное закрытие/\allowbreak{}нет комментария; 42501: нет прав.
\item Гарантии: Фактический конец, комментарий и completed сохраняются атомарно.
\item Трассировка: UC-01; FR-19,20; BR-8,9; T32–34,52,53,57.
\end{itemize}
\subsection{create\_\allowbreak{}reminders}
Создаёт напоминания о будущих подтверждённых бронях до p\_\allowbreak{}until; запускается планировщиком.
\begin{itemize}
\item Вход: p\_\allowbreak{}until timestamptz.
\item Выход: integer: число новых уведомлений.
\item Ошибки: Нет ожидаемых бизнес-ошибок.
\item Гарантии: Частичный UNIQUE и ON CONFLICT: идемпотентность.
\item Трассировка: UC-01; FR-21; T43.
\end{itemize}
\subsection{equipment\_\allowbreak{}metrics}
Плановая и фактическая загрузка, объединённый простой, процент утилизации и число отказов за окно.
\begin{itemize}
\item Вход: p\_\allowbreak{}window tstzrange.
\item Выход: equipment\_\allowbreak{}id, planned\_\allowbreak{}seconds, actual\_\allowbreak{}seconds, downtime\_\allowbreak{}seconds, utilization\_\allowbreak{}percent, conflicts.
\item Ошибки: 22023: окно не конечное/\allowbreak{}пустое.
\item Гарантии: Сначала union диапазонов, затем пересечение с окном; без двойного счёта.
\item Трассировка: UC-03; FR-22,23; T39,56,57.
\end{itemize}
\subsection{grant\_\allowbreak{}admission}
Администратор назначает допуск с основанием и сроком.
\begin{itemize}
\item Вход: p\_\allowbreak{}actor bigint,p\_\allowbreak{}user bigint,p\_\allowbreak{}type bigint,p\_\allowbreak{}start timestamptz,  p\_\allowbreak{}end timestamptz,p\_\allowbreak{}basis text.
\item Выход: bigint: id допуска.
\item Ошибки: 42501: нет роли admin; 23514: интервал; 23503: ссылки.
\item Гарантии: Допуск и аудит сохраняются в транзакции.
\item Трассировка: UC-03; FR-04; BR-1; T45–48.
\end{itemize}
\subsection{revoke\_\allowbreak{}admission}
Администратор отзывает допуск без удаления истории.
\begin{itemize}
\item Вход: p\_\allowbreak{}actor bigint,p\_\allowbreak{}admission bigint.
\item Выход: void.
\item Ошибки: 42501: роль; 23514: уже отозван/\allowbreak{}не найден.
\item Гарантии: При повторном подтверждении и начале работы отзыв учитывается.
\item Трассировка: UC-03; FR-04; BR-1,8; T49,50.
\end{itemize}

\section{Индексы и обоснование}

Созданы 13 явных индексов. PK/UNIQUE/EXCLUDE автоматически создают ещё 25,
всего 38. B-tree соответствует равенству и сортировке, GiST --- покрытию и
пересечению временных диапазонов. Частичные индексы уменьшают объём для нужных
статусов. FK сами по себе индекс дочернего столбца не создают.
Все индексы требуют места, WAL и обслуживания при записи.
\begin{singlespace}\small
\begin{longtable}{>{\raggedright\arraybackslash}p{0.3\textwidth}>{\raggedright\arraybackslash}p{0.31\textwidth}>{\raggedright\arraybackslash}p{0.31\textwidth}}
\toprule Индекс и тип & Ожидаемый запрос & Цена и ограничения\\\midrule\endfirsthead
\toprule Индекс и тип & Ожидаемый запрос & Цена и ограничения\\\midrule\endhead
ix\_\allowbreak{}equipment\_\allowbreak{}type\_\allowbreak{}active; B-tree, частичный & Фильтр активного каталога по типу; компактнее полного индекса. & Архивирование меняет состав; для всех строк не используется. \\
ix\_\allowbreak{}type\_\allowbreak{}capability\_\allowbreak{}reverse; B-tree, составной & Поиск типов по функции; PK начинается с type\_\allowbreak{}id. & Дополнительная запись при изменении M:N. \\
ix\_\allowbreak{}admission\_\allowbreak{}valid; GiST, частичный & Пользователь, тип и покрытие временного диапазона. & GiST дороже B-tree; отзыв обновляет индекс; на 12 строках Seq Scan дешевле. \\
ix\_\allowbreak{}booking\_\allowbreak{}user\_\allowbreak{}future; B-tree, частичный & Календарь владельца, нижняя граница и LIMIT по времени. & Предикат по статусу, не по now(); завершение удаляет индексную запись. \\
ix\_\allowbreak{}booking\_\allowbreak{}equipment\_\allowbreak{}period; GiST, полный & Историческая аналитика и периоды всех статусов. & Для confirmed дублирует часть возможностей EXCLUDE, но нужен для истории. \\
ix\_\allowbreak{}unavailability\_\allowbreak{}period; GiST, частичный & Пересечение ТО/\allowbreak{}поломок с запросом и состояние на момент времени. & Цена записи диапазонов; исключает только отменённые причины, включает resolved. \\
ix\_\allowbreak{}alternative\_\allowbreak{}booking\_\allowbreak{}rank; B-tree, частичный & Список предложений по броне, приоритету и времени. & После выбора требуется менять индексный состав; варианты не резервируют время. \\
ix\_\allowbreak{}fault\_\allowbreak{}equipment\_\allowbreak{}time; B-tree, составной & Последние сообщения по оборудованию. & Занимает память; ускорение не измерено на малом журнале. \\
ix\_\allowbreak{}repair\_\allowbreak{}unavailability; B-tree, составной & История работ по причине недоступности. & Дополнительная стоимость append-only вставки. \\
ix\_\allowbreak{}audit\_\allowbreak{}record\_\allowbreak{}time; B-tree, выражение & История объекта по имени таблицы, id и времени. & JSONB ключ вычисляется на вставке; для связующих пар id NULL, нужен иной запрос. \\
ix\_\allowbreak{}notification\_\allowbreak{}unread; B-tree, частичный & Непрочитанные уведомления владельца. & Прочтение удаляет индексную запись. \\
ux\_\allowbreak{}notification\_\allowbreak{}reminder; B-tree, UNIQUE, частичный & Идемпотентный запуск планировщика. & Одна запись reminder на бронь; после переноса повторная политика требует уточнения. \\
ix\_\allowbreak{}conflict\_\allowbreak{}equipment\_\allowbreak{}time; B-tree, составной & Отказы по оборудованию и времени события. & Статистика отражает только API request\_\allowbreak{}booking. \\
\bottomrule\end{longtable}\end{singlespace}

Полные SQL-определения приведены в \code{sql/indexes/04_indexes.sql} и
\code{docs/INDEXES.md}. EXCLUDE также ускоряет поиск confirmed-overlap.
Полный GiST бронирований предназначен для запросов истории всех статусов;
для confirmed его возможности частично перекрываются с EXCLUDE.

\subsection{Проверка планов запросов}
В тестовой транзакции добавлены 1\,000 единиц и 20\,000 непересекающихся броней,
выполнен ANALYZE, затем EXPLAIN (ANALYZE, BUFFERS). После измерения данные
откатаны. Планы и буферы сохранены в \code{docs/explain-results.txt}.

\begin{center}\small
\begin{tabularx}{\textwidth}{p{.1\textwidth}Xr}
\toprule Запрос & План первого сохранённого полного прогона & Время, мс \\
\midrule
Q1 & Index Scan ex\_booking\_overlap & 0.058 \\
Q2 & Index Scan ix\_booking\_user\_future & 0.019 \\
Q3 & equipment\_pkey и ix\_type\_capability\_reverse & 0.080 \\
Q4 & Seq Scan по 12 допускам & 0.009 \\
Q5 & Index Scan ix\_audit\_record\_time & 0.010 \\
Q1' & После удаления полного GiST остаётся EXCLUDE & 0.031 \\
\bottomrule
\end{tabularx}\end{center}

Замеры единичные, на локальном прогретом кэше. Удаление полного GiST не ухудшило
Q1: оба раза использован индекс EXCLUDE. Это не доказательство ускорения
относительно БД без индексов. На 12 допусках Seq Scan ожидаемо дешевле.
NFR-01 и NFR-02 для веб-сессий и целевого Helios не проверены.

\section{Тестирование реализации}

Фактическое окружение: Fedora 44, PostgreSQL 18.6, локальный кластер в /tmp,
Unix-сокет /tmp, порт 55432, TCP отключён. Использовалась только
\code{labreserve_stage02_test}. Полный runner успешно выполнил создание,
загрузку, проверки, планы, очистку учебной схемы и повторное создание с нуля.

Выполнено \textbf{60 из 60 SQL-утверждений} и \textbf{5 из 5 конкурентных сценариев}.
Ожидаемые запрещённые действия проверяют SQLSTATE; отсутствие ошибки считается
провалом теста. SQL-тесты и временная роль NOLOGIN откатываются одной транзакцией.
\begin{singlespace}\small
\begin{longtable}{>{\raggedright\arraybackslash}p{0.1\textwidth}>{\raggedright\arraybackslash}p{0.68\textwidth}>{\raggedright\arraybackslash}p{0.13\textwidth}}
\toprule Тест & Проверяемое поведение & Результат\\\midrule\endfirsthead
\toprule Тест & Проверяемое поведение & Результат\\\midrule\endhead
T01 & Корректная бронь & PASS \\
T02 & Отказ без допуска & PASS \\
T03 & Неактивный пользователь, прямой DML & PASS \\
T04 & Пересечение броней & PASS \\
T05 & Смежные полуоткрытые интервалы & PASS \\
T06 & Отмена освобождает слот & PASS \\
T07 & Отменённая бронь неизменяема & PASS \\
T08 & Пустой интервал отклонён & PASS \\
T09 & Неподдерживаемая функция & PASS \\
T10 & Открытие до планового начала отклонено & PASS \\
T11 & Проверка владельца & PASS \\
T12 & ТО переводит бронь в needs\_\allowbreak{}transfer & PASS \\
T13 & Не менее трёх альтернатив & PASS \\
T14 & Атомарное уведомление & PASS \\
T15 & Бронь во время ТО отклонена & PASS \\
T16 & Выбор альтернативы & PASS \\
T17 & Выбор сохранён в истории & PASS \\
T18 & Повторный выбор отклонён & PASS \\
T19 & Приоритет исходного оборудования BR-6 & PASS \\
T20 & Отмена ТО не подтверждает брони & PASS \\
T21 & Сообщение сохраняется без блокировки & PASS \\
T22 & Администратор подтверждает причину & PASS \\
T23 & Поломка затрагивает будущую бронь & PASS \\
T24 & Неисправное оборудование недоступно & PASS \\
T25 & Пользователь не подтверждает причину & PASS \\
T26 & Поломка без конца исключает исходный станок & PASS \\
T27 & Другие причины сохраняют недоступность & PASS \\
T28 & Отсутствие альтернатив & PASS \\
T29 & Техник восстанавливает доступность & PASS \\
T30 & Корректное открытие сеанса & PASS \\
T31 & Один сеанс на бронь & PASS \\
T32 & Закрытие завершает бронь & PASS \\
T33 & Повторное закрытие отклонено & PASS \\
T34 & Завершённая бронь неизменяема & PASS \\
T35 & Аудит неизменяем & PASS \\
T36 & Значимые изменения отражены в аудите & PASS \\
T37 & Равные моменты разных зон & PASS \\
T38 & Отклонённый запрос сохранён & PASS \\
T39 & Ограничение метрик окном & PASS \\
T40 & Эквивалентность разных типов & PASS \\
T41 & Удаление брони запрещено & PASS \\
T42 & Поиск максимальных свободных интервалов & PASS \\
T43 & Напоминания идемпотентны & PASS \\
T44 & Конкурентно занятая альтернатива отклонена & PASS \\
T45 & Допуск назначает администратор & PASS \\
T46 & Поиск учитывает будущий срок допуска & PASS \\
T47 & Не наступивший допуск отклонён & PASS \\
T48 & Допуск покрывает весь интервал & PASS \\
T49 & Отозванный допуск отклонён & PASS \\
T50 & История допуска сохраняется & PASS \\
T51 & История недоступности сохраняется & PASS \\
T52 & Администратору нужен комментарий закрытия & PASS \\
T53 & Комментарий закрытия сохранён & PASS \\
T54 & Включённый конец диапазона отклонён & PASS \\
T55 & В схеме 17 таблиц & PASS \\
T56 & Перекрывающиеся простои считаются объединением & PASS \\
T57 & Фактическая длительность отделена от плановой & PASS \\
T58 & Backend не имеет прямого DML расписания & PASS \\
T59 & Backend не вызывает внутренний guard & PASS \\
T60 & API работает без session search\_\allowbreak{}path & PASS \\
\bottomrule\end{longtable}\end{singlespace}
\begin{singlespace}\small
\begin{longtable}{>{\raggedright\arraybackslash}p{0.1\textwidth}>{\raggedright\arraybackslash}p{0.68\textwidth}>{\raggedright\arraybackslash}p{0.13\textwidth}}
\toprule Тест & Проверяемое поведение & Результат\\\midrule\endfirsthead
\toprule Тест & Проверяемое поведение & Результат\\\midrule\endhead
N01 & вид работы при устранении неисправности & PASS \\
N02 & вид работы при завершении ТО & PASS \\
N03 & пустой вид работы отклонён & PASS \\
N04 & вид работы обязателен & PASS \\
\bottomrule\end{longtable}\end{singlespace}

\subsection{Два конкурентных соединения}
Python-скрипт использует стандартную библиотеку и два процесса psql.
Первая транзакция удерживает блокировку; вторая запускается только после
наблюдения PgSleep первой в pg\_stat\_activity. Проверяются код отказа и
итоговое отсутствие подтверждённого пересечения.
\begin{enumerate}
\item C01: две одинаковые брони; первая зафиксирована, вторая получает 23P01,
подтверждена одна бронь.
\item C02: ТО фиксируется раньше брони; бронь получает 23514.
\item C03: бронь фиксируется раньше ТО; ТО успешно, бронь получает needs\_transfer.
\item C04: два прямых INSERT; второй получает 23P01.
\item C05: REPEATABLE READ со снимком до COMMIT ТО; второй UPDATE оборудования
получает 40001, устаревшая проверка не выполняется.
\end{enumerate}
Результаты с серверными сообщениями сохранены в \code{docs/concurrency-results.json}.

\subsection{Препятствия и исправления}
\begin{enumerate}
\item Песочница запрещала Unix-сокет. Разрешённый запуск вне песочницы позволил
проверить изолированный кластер без обращения к пользовательским БД.
\item В системе отсутствовал contrib. Совпадающий RPM PostgreSQL 18.6
загружен и извлечён в /tmp, подключён к временной конфигурации; системные
библиотеки PostgreSQL не заменялись.
\item При первой установке исправлена пропущенная скобка домена; транзакция
откатана и установка выполнена заново.
\item В T30 исправлен тестовый интервал, который ошибочно захватывал закрытую
поломку. База правильно отклоняла этот запрос.
\item В T38 тест отказа перенесён с уже восстановленного оборудования на
действительно неисправное.
\item При визуальной проверке диаграммы переставлены узлы и вынесены подписи
FK, которые в первой версии перекрывали текст таблиц.
\end{enumerate}
Первый этап проверяется сравнением SHA-256 всех файлов до и после работы.
История тестов и ограничения окружения изложены в \code{docs/TESTING.md}.

\section{Рабочие решения и вопросы согласования}

\begin{enumerate}
\item Подтвердить ER-модель из 15 сущностей, две M:N и единый период
недоступности для ТО и подтверждённой неисправности.
\item Подтвердить допуск на весь интервал. Отзыв не переписывает историю,
но запрещает новые подтверждения и начало работы; массовая обработка будущих
броней после отзыва пока требует политики.
\item Уточнить эквивалентность: достаточно одной функции либо нужны
материал, мощность, точность и набор функций.
\item Подтвердить поиск вперёд, горизонт 30 суток и представление слотов:
в каждом свободном интервале перечисляются последовательные непересекающиеся
работы исходной длительности. Это не перебор всех непрерывных начал.
\item Уточнить остановку уже начатой работы при поломке. BR-7 касается будущих
броней; автоматическое завершение начатого сеанса не выполняется.
\item Подтвердить календарный знаменатель утилизации. Рабочий календарь,
сменность и «простой без заявок» пока не определены.
\item Уточнить повторные напоминания после переноса и массовую обработку
назначенных броней при архивировании каталога.
\item Проверить версии PostgreSQL и право установки btree\_gist на Helios.
Локально проверена 18.6; использованные multirange требуют PostgreSQL 14+.
\end{enumerate}

Не выполнены внешнее согласование ER, развёртывание на Helios и нагрузочная
проверка всего веб-приложения. Эти границы не препятствуют проверке базы второго
этапа. Статус approved к модели не применяется.

\section{Заключение}

Созданы ER- и даталогическая модели, PostgreSQL-схема из 17 таблиц,
SQL API критических операций, индексы и проверочный комплект.
Все 60 SQL-утверждений и пять конкурентных сценариев прошли.
Материалы пригодны для демонстрации второго этапа и последующей интеграции
через JPA/Hibernate. ER-модель ожидает согласования преподавателем.

\clearpage
\appendix
\section{Полный словарь таблиц, ключей и ограничений}
\label{app:dictionary}
Все поля, обязательность, default и определения ограничений ниже получены из
реально установленной PostgreSQL-схемы. FK везде используют RESTRICT при
удалении и обновлении. На диаграммах показан базовый тип tstzrange; в DDL
использован проверяемый домен time\_window. Первичные ключи IDENTITY генерируются
СУБД. Индексы описаны в основном разделе и отдельном SQL-файле.
\begingroup\small\setstretch{1.1}
\par\noindent\begin{minipage}{\linewidth}
\subsection{app\_\allowbreak{}user}
Учётная запись человека; активность и профиль.\par\smallskip\noindent
\begin{tabular}{p{0.29\textwidth}p{0.25\textwidth}p{0.1\textwidth}p{0.25\textwidth}}
\toprule Столбец & Тип & NULL & По умолчанию \\ \midrule
id & bigint & нет & IDENTITY \\
login & text & нет & — \\
full\_\allowbreak{}name & text & нет & — \\
password\_\allowbreak{}hash & text & да & — \\
active & boolean & нет & true \\
created\_\allowbreak{}at & timestamptz & нет & clock\_\allowbreak{}timestamp() \\
\bottomrule\end{tabular}\end{minipage}\par
\begin{itemize}
\item CHECK ((length(TRIM(BOTH FROM login)) > 0)).
\item PRIMARY KEY (id).
\item UNIQUE (login).
\end{itemize}
\par\noindent\begin{minipage}{\linewidth}
\subsection{app\_\allowbreak{}role}
Три внешние роли предметной области.\par\smallskip\noindent
\begin{tabular}{p{0.29\textwidth}p{0.25\textwidth}p{0.1\textwidth}p{0.25\textwidth}}
\toprule Столбец & Тип & NULL & По умолчанию \\ \midrule
code & text & нет & — \\
\bottomrule\end{tabular}\end{minipage}\par
\begin{itemize}
\item CHECK ((code = ANY (ARRAY['user'::text, 'admin'::text, 'technician'::text]))).
\item PRIMARY KEY (code).
\end{itemize}
\par\noindent\begin{minipage}{\linewidth}
\subsection{user\_\allowbreak{}role}
Назначение ролей пользователям (M:N).\par\smallskip\noindent
\begin{tabular}{p{0.29\textwidth}p{0.25\textwidth}p{0.1\textwidth}p{0.25\textwidth}}
\toprule Столбец & Тип & NULL & По умолчанию \\ \midrule
user\_\allowbreak{}id & bigint & нет & — \\
role\_\allowbreak{}code & text & нет & — \\
\bottomrule\end{tabular}\end{minipage}\par
\begin{itemize}
\item PRIMARY KEY (user\_\allowbreak{}id, role\_\allowbreak{}code).
\item FOREIGN KEY (user\_\allowbreak{}id) REFERENCES labreserve.app\_\allowbreak{}user(id) ON UPDATE RESTRICT ON DELETE RESTRICT.
\item FOREIGN KEY (role\_\allowbreak{}code) REFERENCES labreserve.app\_\allowbreak{}role(code) ON UPDATE RESTRICT ON DELETE RESTRICT.
\end{itemize}
\par\noindent\begin{minipage}{\linewidth}
\subsection{equipment\_\allowbreak{}type}
Тип оборудования и архивирование.\par\smallskip\noindent
\begin{tabular}{p{0.29\textwidth}p{0.25\textwidth}p{0.1\textwidth}p{0.25\textwidth}}
\toprule Столбец & Тип & NULL & По умолчанию \\ \midrule
id & bigint & нет & IDENTITY \\
name & text & нет & — \\
archived & boolean & нет & false \\
\bottomrule\end{tabular}\end{minipage}\par
\begin{itemize}
\item PRIMARY KEY (id).
\item UNIQUE (name).
\end{itemize}
\par\noindent\begin{minipage}{\linewidth}
\subsection{capability}
Справочник требуемых функций.\par\smallskip\noindent
\begin{tabular}{p{0.29\textwidth}p{0.25\textwidth}p{0.1\textwidth}p{0.25\textwidth}}
\toprule Столбец & Тип & NULL & По умолчанию \\ \midrule
id & bigint & нет & IDENTITY \\
name & text & нет & — \\
\bottomrule\end{tabular}\end{minipage}\par
\begin{itemize}
\item PRIMARY KEY (id).
\item UNIQUE (name).
\end{itemize}
\par\noindent\begin{minipage}{\linewidth}
\subsection{type\_\allowbreak{}capability}
Поддерживаемые типом функции (M:N).\par\smallskip\noindent
\begin{tabular}{p{0.29\textwidth}p{0.25\textwidth}p{0.1\textwidth}p{0.25\textwidth}}
\toprule Столбец & Тип & NULL & По умолчанию \\ \midrule
type\_\allowbreak{}id & bigint & нет & — \\
capability\_\allowbreak{}id & bigint & нет & — \\
\bottomrule\end{tabular}\end{minipage}\par
\begin{itemize}
\item PRIMARY KEY (type\_\allowbreak{}id, capability\_\allowbreak{}id).
\item FOREIGN KEY (type\_\allowbreak{}id) REFERENCES labreserve.equipment\_\allowbreak{}type(id) ON UPDATE RESTRICT ON DELETE RESTRICT.
\item FOREIGN KEY (capability\_\allowbreak{}id) REFERENCES labreserve.capability(id) ON UPDATE RESTRICT ON DELETE RESTRICT.
\end{itemize}
\par\noindent\begin{minipage}{\linewidth}
\subsection{equipment}
Инвентарная единица, характеристики, архив.\par\smallskip\noindent
\begin{tabular}{p{0.29\textwidth}p{0.25\textwidth}p{0.1\textwidth}p{0.25\textwidth}}
\toprule Столбец & Тип & NULL & По умолчанию \\ \midrule
id & bigint & нет & IDENTITY \\
type\_\allowbreak{}id & bigint & нет & — \\
inventory\_\allowbreak{}no & text & нет & — \\
name & text & нет & — \\
characteristics & jsonb & нет & '\{\}'::jsonb \\
archived & boolean & нет & false \\
schedule\_\allowbreak{}revision & bigint & нет & 0 \\
\bottomrule\end{tabular}\end{minipage}\par
\begin{itemize}
\item CHECK ((jsonb\_\allowbreak{}typeof(characteristics) = 'object'::text)).
\item PRIMARY KEY (id).
\item UNIQUE (inventory\_\allowbreak{}no).
\item FOREIGN KEY (type\_\allowbreak{}id) REFERENCES labreserve.equipment\_\allowbreak{}type(id) ON UPDATE RESTRICT ON DELETE RESTRICT.
\end{itemize}
\par\noindent\begin{minipage}{\linewidth}
\subsection{admission}
Допуск к типу; период, основание, автор и отзыв.\par\smallskip\noindent
\begin{tabular}{p{0.29\textwidth}p{0.25\textwidth}p{0.1\textwidth}p{0.25\textwidth}}
\toprule Столбец & Тип & NULL & По умолчанию \\ \midrule
id & bigint & нет & IDENTITY \\
user\_\allowbreak{}id & bigint & нет & — \\
type\_\allowbreak{}id & bigint & нет & — \\
valid\_\allowbreak{}during & time\_\allowbreak{}window & нет & — \\
basis & text & нет & — \\
granted\_\allowbreak{}by & bigint & нет & — \\
revoked\_\allowbreak{}at & timestamptz & да & — \\
created\_\allowbreak{}at & timestamptz & нет & clock\_\allowbreak{}timestamp() \\
\bottomrule\end{tabular}\end{minipage}\par
\begin{itemize}
\item PRIMARY KEY (id).
\item FOREIGN KEY (user\_\allowbreak{}id) REFERENCES labreserve.app\_\allowbreak{}user(id) ON UPDATE RESTRICT ON DELETE RESTRICT.
\item FOREIGN KEY (type\_\allowbreak{}id) REFERENCES labreserve.equipment\_\allowbreak{}type(id) ON UPDATE RESTRICT ON DELETE RESTRICT.
\item FOREIGN KEY (granted\_\allowbreak{}by) REFERENCES labreserve.app\_\allowbreak{}user(id) ON UPDATE RESTRICT ON DELETE RESTRICT.
\end{itemize}
\par\noindent\begin{minipage}{\linewidth}
\subsection{booking}
Плановый интервал, владелец, цель и жизненный цикл.\par\smallskip\noindent
\begin{tabular}{p{0.29\textwidth}p{0.25\textwidth}p{0.1\textwidth}p{0.25\textwidth}}
\toprule Столбец & Тип & NULL & По умолчанию \\ \midrule
id & bigint & нет & IDENTITY \\
user\_\allowbreak{}id & bigint & нет & — \\
equipment\_\allowbreak{}id & bigint & нет & — \\
capability\_\allowbreak{}id & bigint & нет & — \\
period & time\_\allowbreak{}window & нет & — \\
purpose & text & нет & — \\
status & text & нет & 'confirmed'::text \\
cancelled\_\allowbreak{}at & timestamptz & да & — \\
cancel\_\allowbreak{}reason & text & да & — \\
created\_\allowbreak{}at & timestamptz & нет & clock\_\allowbreak{}timestamp() \\
\bottomrule\end{tabular}\end{minipage}\par
\begin{itemize}
\item CHECK ((NOT upper\_\allowbreak{}inf((period)::tstzrange))).
\item CHECK ((status = ANY (ARRAY['confirmed'::text, 'needs\_\allowbreak{}transfer'::text, 'cancelled'::text, 'completed'::text]))).
\item CHECK (((status = 'cancelled'::text) = ((cancelled\_\allowbreak{}at IS NOT NULL) AND (cancel\_\allowbreak{}reason IS NOT NULL)))).
\item PRIMARY KEY (id).
\item EXCLUDE USING gist (equipment\_\allowbreak{}id WITH =, ((period)::tstzrange) WITH \&\&) WHERE ((status = 'confirmed'::text)).
\item FOREIGN KEY (user\_\allowbreak{}id) REFERENCES labreserve.app\_\allowbreak{}user(id) ON UPDATE RESTRICT ON DELETE RESTRICT.
\item FOREIGN KEY (equipment\_\allowbreak{}id) REFERENCES labreserve.equipment(id) ON UPDATE RESTRICT ON DELETE RESTRICT.
\item FOREIGN KEY (capability\_\allowbreak{}id) REFERENCES labreserve.capability(id) ON UPDATE RESTRICT ON DELETE RESTRICT.
\end{itemize}
\par\noindent\begin{minipage}{\linewidth}
\subsection{alternative}
Предложение переноса: интервал, приоритет и статус.\par\smallskip\noindent
\begin{tabular}{p{0.29\textwidth}p{0.25\textwidth}p{0.1\textwidth}p{0.25\textwidth}}
\toprule Столбец & Тип & NULL & По умолчанию \\ \midrule
id & bigint & нет & IDENTITY \\
booking\_\allowbreak{}id & bigint & нет & — \\
equipment\_\allowbreak{}id & bigint & нет & — \\
period & time\_\allowbreak{}window & нет & — \\
priority & smallint & нет & — \\
state & text & нет & 'offered'::text \\
created\_\allowbreak{}at & timestamptz & нет & clock\_\allowbreak{}timestamp() \\
\bottomrule\end{tabular}\end{minipage}\par
\begin{itemize}
\item CHECK ((NOT upper\_\allowbreak{}inf((period)::tstzrange))).
\item CHECK (((priority >= 1) AND (priority <= 3))).
\item CHECK ((state = ANY (ARRAY['offered'::text, 'chosen'::text, 'obsolete'::text]))).
\item PRIMARY KEY (id).
\item UNIQUE (booking\_\allowbreak{}id, equipment\_\allowbreak{}id, period).
\item FOREIGN KEY (booking\_\allowbreak{}id) REFERENCES labreserve.booking(id) ON UPDATE RESTRICT ON DELETE RESTRICT.
\item FOREIGN KEY (equipment\_\allowbreak{}id) REFERENCES labreserve.equipment(id) ON UPDATE RESTRICT ON DELETE RESTRICT.
\end{itemize}
\par\noindent\begin{minipage}{\linewidth}
\subsection{fault\_\allowbreak{}report}
Первичное сообщение и результат проверки.\par\smallskip\noindent
\begin{tabular}{p{0.29\textwidth}p{0.25\textwidth}p{0.1\textwidth}p{0.25\textwidth}}
\toprule Столбец & Тип & NULL & По умолчанию \\ \midrule
id & bigint & нет & IDENTITY \\
equipment\_\allowbreak{}id & bigint & нет & — \\
reported\_\allowbreak{}by & bigint & нет & — \\
description & text & нет & — \\
reported\_\allowbreak{}at & timestamptz & нет & clock\_\allowbreak{}timestamp() \\
state & text & нет & 'reported'::text \\
\bottomrule\end{tabular}\end{minipage}\par
\begin{itemize}
\item CHECK ((state = ANY (ARRAY['reported'::text, 'confirmed'::text, 'dismissed'::text]))).
\item PRIMARY KEY (id).
\item FOREIGN KEY (equipment\_\allowbreak{}id) REFERENCES labreserve.equipment(id) ON UPDATE RESTRICT ON DELETE RESTRICT.
\item FOREIGN KEY (reported\_\allowbreak{}by) REFERENCES labreserve.app\_\allowbreak{}user(id) ON UPDATE RESTRICT ON DELETE RESTRICT.
\end{itemize}
\par\noindent\begin{minipage}{\linewidth}
\subsection{unavailability}
Техническое событие: ТО либо подтверждённая неисправность.\par\smallskip\noindent
\begin{tabular}{p{0.29\textwidth}p{0.25\textwidth}p{0.1\textwidth}p{0.25\textwidth}}
\toprule Столбец & Тип & NULL & По умолчанию \\ \midrule
id & bigint & нет & IDENTITY \\
equipment\_\allowbreak{}id & bigint & нет & — \\
kind & text & нет & — \\
period & time\_\allowbreak{}window & нет & — \\
reason & text & нет & — \\
fault\_\allowbreak{}report\_\allowbreak{}id & bigint & да & — \\
created\_\allowbreak{}by & bigint & нет & — \\
status & text & нет & 'active'::text \\
resolved\_\allowbreak{}by & bigint & да & — \\
resolved\_\allowbreak{}at & timestamptz & да & — \\
created\_\allowbreak{}at & timestamptz & нет & clock\_\allowbreak{}timestamp() \\
\bottomrule\end{tabular}\end{minipage}\par
\begin{itemize}
\item CHECK ((kind = ANY (ARRAY['maintenance'::text, 'fault'::text]))).
\item CHECK ((status = ANY (ARRAY['active'::text, 'resolved'::text, 'cancelled'::text]))).
\item CHECK (((kind = 'fault'::text) = (fault\_\allowbreak{}report\_\allowbreak{}id IS NOT NULL))).
\item CHECK (((kind <> 'maintenance'::text) OR (NOT upper\_\allowbreak{}inf((period)::tstzrange)))).
\item CHECK (((status = 'resolved'::text) = ((resolved\_\allowbreak{}by IS NOT NULL) AND (resolved\_\allowbreak{}at IS NOT NULL)))).
\item PRIMARY KEY (id).
\item UNIQUE (fault\_\allowbreak{}report\_\allowbreak{}id).
\item FOREIGN KEY (equipment\_\allowbreak{}id) REFERENCES labreserve.equipment(id) ON UPDATE RESTRICT ON DELETE RESTRICT.
\item FOREIGN KEY (fault\_\allowbreak{}report\_\allowbreak{}id) REFERENCES labreserve.fault\_\allowbreak{}report(id) ON UPDATE RESTRICT ON DELETE RESTRICT.
\item FOREIGN KEY (created\_\allowbreak{}by) REFERENCES labreserve.app\_\allowbreak{}user(id) ON UPDATE RESTRICT ON DELETE RESTRICT.
\item FOREIGN KEY (resolved\_\allowbreak{}by) REFERENCES labreserve.app\_\allowbreak{}user(id) ON UPDATE RESTRICT ON DELETE RESTRICT.
\end{itemize}
\par\noindent\begin{minipage}{\linewidth}
\subsection{repair}
Техническая работа: вид, исполнитель, интервал и результат.\par\smallskip\noindent
\begin{tabular}{p{0.29\textwidth}p{0.25\textwidth}p{0.1\textwidth}p{0.25\textwidth}}
\toprule Столбец & Тип & NULL & По умолчанию \\ \midrule
id & bigint & нет & IDENTITY \\
unavailability\_\allowbreak{}id & bigint & нет & — \\
performed\_\allowbreak{}by & bigint & нет & — \\
work\_\allowbreak{}kind & text & нет & — \\
period & time\_\allowbreak{}window & нет & — \\
result & text & нет & — \\
\bottomrule\end{tabular}\end{minipage}\par
\begin{itemize}
\item CHECK ((length(TRIM(BOTH FROM work\_\allowbreak{}kind)) > 0)).
\item CHECK ((NOT upper\_\allowbreak{}inf((period)::tstzrange))).
\item PRIMARY KEY (id).
\item FOREIGN KEY (unavailability\_\allowbreak{}id) REFERENCES labreserve.unavailability(id) ON UPDATE RESTRICT ON DELETE RESTRICT.
\item FOREIGN KEY (performed\_\allowbreak{}by) REFERENCES labreserve.app\_\allowbreak{}user(id) ON UPDATE RESTRICT ON DELETE RESTRICT.
\end{itemize}
\par\noindent\begin{minipage}{\linewidth}
\subsection{use\_\allowbreak{}session}
Фактические начало и окончание работы.\par\smallskip\noindent
\begin{tabular}{p{0.29\textwidth}p{0.25\textwidth}p{0.1\textwidth}p{0.25\textwidth}}
\toprule Столбец & Тип & NULL & По умолчанию \\ \midrule
id & bigint & нет & IDENTITY \\
booking\_\allowbreak{}id & bigint & нет & — \\
started\_\allowbreak{}at & timestamptz & нет & clock\_\allowbreak{}timestamp() \\
ended\_\allowbreak{}at & timestamptz & да & — \\
close\_\allowbreak{}comment & text & да & — \\
\bottomrule\end{tabular}\end{minipage}\par
\begin{itemize}
\item CHECK (((ended\_\allowbreak{}at IS NULL) OR (ended\_\allowbreak{}at > started\_\allowbreak{}at))).
\item PRIMARY KEY (id).
\item UNIQUE (booking\_\allowbreak{}id).
\item FOREIGN KEY (booking\_\allowbreak{}id) REFERENCES labreserve.booking(id) ON UPDATE RESTRICT ON DELETE RESTRICT.
\end{itemize}
\par\noindent\begin{minipage}{\linewidth}
\subsection{notification}
Внутрисистемное уведомление и прочтение.\par\smallskip\noindent
\begin{tabular}{p{0.29\textwidth}p{0.25\textwidth}p{0.1\textwidth}p{0.25\textwidth}}
\toprule Столбец & Тип & NULL & По умолчанию \\ \midrule
id & bigint & нет & IDENTITY \\
user\_\allowbreak{}id & bigint & нет & — \\
booking\_\allowbreak{}id & bigint & да & — \\
kind & text & нет & — \\
message & text & нет & — \\
created\_\allowbreak{}at & timestamptz & нет & clock\_\allowbreak{}timestamp() \\
read\_\allowbreak{}at & timestamptz & да & — \\
\bottomrule\end{tabular}\end{minipage}\par
\begin{itemize}
\item CHECK ((kind = ANY (ARRAY['transfer\_\allowbreak{}required'::text, 'transferred'::text, 'cancelled'::text, 'reminder'::text]))).
\item PRIMARY KEY (id).
\item FOREIGN KEY (user\_\allowbreak{}id) REFERENCES labreserve.app\_\allowbreak{}user(id) ON UPDATE RESTRICT ON DELETE RESTRICT.
\item FOREIGN KEY (booking\_\allowbreak{}id) REFERENCES labreserve.booking(id) ON UPDATE RESTRICT ON DELETE RESTRICT.
\end{itemize}
\par\noindent\begin{minipage}{\linewidth}
\subsection{audit\_\allowbreak{}event}
Служебный журнал значимых изменений (вне инфологической модели).\par\smallskip\noindent
\begin{tabular}{p{0.29\textwidth}p{0.25\textwidth}p{0.1\textwidth}p{0.25\textwidth}}
\toprule Столбец & Тип & NULL & По умолчанию \\ \midrule
id & bigint & нет & IDENTITY \\
actor\_\allowbreak{}id & bigint & да & — \\
table\_\allowbreak{}name & text & нет & — \\
record\_\allowbreak{}key & jsonb & нет & — \\
action & text & нет & — \\
old\_\allowbreak{}data & jsonb & да & — \\
new\_\allowbreak{}data & jsonb & да & — \\
occurred\_\allowbreak{}at & timestamptz & нет & clock\_\allowbreak{}timestamp() \\
transaction\_\allowbreak{}id & bigint & нет & txid\_\allowbreak{}current() \\
\bottomrule\end{tabular}\end{minipage}\par
\begin{itemize}
\item PRIMARY KEY (id).
\item FOREIGN KEY (actor\_\allowbreak{}id) REFERENCES labreserve.app\_\allowbreak{}user(id) ON UPDATE RESTRICT ON DELETE RESTRICT.
\end{itemize}
\par\noindent\begin{minipage}{\linewidth}
\subsection{conflict\_\allowbreak{}event}
Служебный журнал отказов для статистики (вне инфологической модели).\par\smallskip\noindent
\begin{tabular}{p{0.29\textwidth}p{0.25\textwidth}p{0.1\textwidth}p{0.25\textwidth}}
\toprule Столбец & Тип & NULL & По умолчанию \\ \midrule
id & bigint & нет & IDENTITY \\
user\_\allowbreak{}id & bigint & нет & — \\
equipment\_\allowbreak{}id & bigint & нет & — \\
requested\_\allowbreak{}period & time\_\allowbreak{}window & нет & — \\
error\_\allowbreak{}code & text & нет & — \\
occurred\_\allowbreak{}at & timestamptz & нет & clock\_\allowbreak{}timestamp() \\
\bottomrule\end{tabular}\end{minipage}\par
\begin{itemize}
\item CHECK ((NOT upper\_\allowbreak{}inf((requested\_\allowbreak{}period)::tstzrange))).
\item PRIMARY KEY (id).
\item FOREIGN KEY (user\_\allowbreak{}id) REFERENCES labreserve.app\_\allowbreak{}user(id) ON UPDATE RESTRICT ON DELETE RESTRICT.
\item FOREIGN KEY (equipment\_\allowbreak{}id) REFERENCES labreserve.equipment(id) ON UPDATE RESTRICT ON DELETE RESTRICT.
\end{itemize}
\endgroup

\clearpage
\section{Примеры исполняемого SQL}
\subsection{Защита подтверждённых интервалов}
Фрагмент определения таблицы booking; остальные столбцы и ограничения находятся
в \code{sql/schema/01_tables.sql}.
\noindent\begin{minipage}{\linewidth}
\begin{Verbatim}[fontsize=\small,breaklines=true]
CONSTRAINT ex_booking_overlap EXCLUDE USING gist
 (equipment_id WITH =, (period::tstzrange) WITH &&)
 WHERE (status='confirmed')
\end{Verbatim}
\end{minipage}

\subsection{Общий ресурс сериализации}
Обновление строки необходимо для отказа на устаревшем снимке более строгой
изоляции; простое чтение или предварительный SELECT не заменяет эту операцию.
\noindent\begin{minipage}{\linewidth}
\begin{Verbatim}[fontsize=\small,breaklines=true]
CREATE FUNCTION lock_equipment(p_id bigint)
RETURNS void LANGUAGE plpgsql AS $$
BEGIN
 UPDATE equipment
 SET schedule_revision=schedule_revision+1 WHERE id=p_id;
 IF NOT FOUND THEN
  RAISE EXCEPTION 'Unknown equipment' USING ERRCODE='23503';
 END IF;
END $$;
\end{Verbatim}
\end{minipage}

\subsection{Воспроизводимый запуск}
Из каталога этапа, после запуска отдельного локального PostgreSQL:
\noindent\begin{minipage}{\linewidth}
\begin{Verbatim}[fontsize=\small,breaklines=true]
export PGHOST=/tmp PGPORT=55432
psql -d postgres -f sql/scripts/create_database.sql
export PGDATABASE=labreserve_stage02_test
bash sql/scripts/run_checks.sh
psql -X -f sql/scripts/export_report.sql > output/analytics.csv
\end{Verbatim}
\end{minipage}

\subsection{Пример вызова следующего этапа}
Параметры native JPA-запроса связываются отдельно, без конкатенации SQL.
Сервер передаёт идентичность аутентифицированного пользователя и выполняет COMMIT
при успехе либо ROLLBACK при исключении.
\noindent\begin{minipage}{\linewidth}
\begin{Verbatim}[fontsize=\small,breaklines=true]
SELECT labreserve.create_booking(
 :user_id, :equipment_id, :capability_id,
 :start_at, :end_at, :purpose);
SELECT labreserve.select_alternative(:actor_id, :alternative_id);
SELECT * FROM labreserve.equipment_metrics(:report_window);
\end{Verbatim}
\end{minipage}

\section{Источники и состав электронного приложения}
\begin{enumerate}
\item Официальное задание: «Информационные системы. Курсовая работа»,
кафедра ВТ, копия в stage-01, раздел «Этап 2».
\item Утверждённый отчёт первого этапа \systemname, stage-01/report\_approved.pdf,
BR-1--BR-10, FR-01--FR-24 и UC-01--UC-03.
\item PostgreSQL 18. Ограничения:
\url{https://www.postgresql.org/docs/18/ddl-constraints.html}.
\item PostgreSQL 18. Типы диапазонов:
\url{https://www.postgresql.org/docs/18/rangetypes.html}.
\item PostgreSQL 18. Изоляция транзакций:
\url{https://www.postgresql.org/docs/18/transaction-iso.html}.
\item PostgreSQL 18. btree\_gist:
\url{https://www.postgresql.org/docs/18/btree-gist.html}.
\end{enumerate}

Электронное приложение в stage-02 содержит SQL всех таблиц, триггеров, операций
и индексов, seed, проверки, скрипты установки/удаления, редактируемые диаграммы,
JSON фактической схемы, контракты и матрицу трассировки. Назначение файлов и
команды запуска приведены в README.md. Журналы и растровые страницы сборки
размещены в output и исключены из Git; отобранные результаты проверок сохранены
в docs. Первый этап не изменялся.
\end{document}
